\section[Initializing the cap induction and removing the auxiliary potential]{Initializing the cap induction and\\removing the auxiliary potential}
\label{cap:initialization}

The induction of Proposition~\ref{cap:curvature-induction} needs a negative
radial correction at its first scale. We obtain this correction from the
small-field TAP expansion, initially for an auxiliary potential. The auxiliary
potential differs from the true cap potential by a bounded Euclidean gradient
error. A finite number of cap steps contracts this error, after which a
covariance comparison permits us to use the true potential at every remaining
scale.

We retain the notation $P_p,k_p,M_p,R_p,a_p$ of the cap construction and
write
\[
 \mathcal A_p(s)=\{h\in P_p:|\alpha_p(h)|\le s p^2\},
 \qquad \phi_p(h)=\log Z_p(h).
\]
Here $\alpha_p(h)$ is determined by $n=\operatorname{Tr}R_p+\|R_ph\|^2$.
Whenever an auxiliary potential $\widetilde\phi_p$ is used, its derivatives
and its discrepancy from $\phi_p$ are restricted to $P_p$. An auxiliary object
is called \emph{intrinsic to the parent data} if it is measurable with respect
to the capped matrix, the eigenvalues, and auxiliary randomness independent of
the as yet unused finer eigenframe. This requirement is essential for the
contraction argument below.

All parameters in the initialization and the finite comparison phase are fixed
before $n$ tends to infinity. Fix first a finite field radius $L_0$ for the
eventual bounded-ball estimate. The construction tolerances may depend on
$L_0$; the coarse constant supplied by
Lemma~\ref{cap:frozen-evaluation-support} depends only on $L_0$ and is
fixed before the sharp curvature tolerances. Choose a sufficiently large absolute $C_*$
and the small excess $\eta_*$ used below, and then the constant $K_0$ in
\eqref{cap:invariant}. Choose the scale ratio
$\gamma$ sufficiently small, the annulus tolerance $\sigma$ a sufficiently
high power of $\gamma$, and $\delta$ a still higher power. Increasing the
exponents is allowed, and the choice of $\delta$ is made last among these
curvature tolerances. The finitely many error tolerances in the static
comparison are chosen after these choices. Finally choose the starting scale
$p_0>0$ sufficiently small. The coefficient $1-\beta>0$ introduced below is
chosen still smaller, including the fixed-scale persistence requirement in
Section~\ref{cap:constant-scales}; its scalar endpoint and buffers are fixed
from $p_0$ before this choice. The coefficient is thereafter fixed. Constants in the companion
fixed-temperature theorem may depend on this $\beta$ and on $p_0$.

\subsection{A terminal TAP potential and its extension}

Given the data at the starting cap $p=p_0$, restore the eigenvalues above the
cap using an independent conditional Haar frame in $H_p$. Conditional on
the capped matrix and eigenvalues, this auxiliary refinement and the
original refinement are independent. Denote the resulting matrix by $W$
and its cap subtraction by $B_p^{\rm aux}$, so that $W-B_p^{\rm aux}=J_p$.
Its unconditional law is GOE. The sigma-field
$\mathscr F_p^{\rm aux}$ generated by the capped data, eigenvalues, and
auxiliary Haar seed excludes the original finer frame. Set $t=C_*p^2$ and observe with
precision
\begin{equation}\label{cap:initial-precision}
 \Delta=tI+(1-\beta)W-B_p^{\rm aux}.
\end{equation}
For $C_*$ large and $1-\beta$ sufficiently small this is positive definite.
The terminal interaction is $\beta W-tI$, which gives the same spin law as
$\beta W_{\rm off}$. Scalar and Gaussian diagonal terms change only the
partition-function normalization.

On the stable terminal region constructed below, let $\Psi(z)$ be the TAP
potential at its unique root, with interaction $J_*=\beta W_{\rm off}$ and
$j=\beta^2$. Thus $\nabla\Psi(z)=m(z)$, the TAP magnetization.
Propositions~\ref{cap:classical-input} and~\ref{cap:tap-expansion} give the
following construction on an event of probability tending to one. We describe
its domain carefully because it need not be convex.

For a starting field $h\in\mathcal A_p(\sigma)$, freeze its parameter
$\alpha=\alpha_p(h)$. Write $\mathsf G_h$ for the canonical Gaussian
terminal law
\[
 Z=h+\Delta a_p+\zeta,\qquad
 \zeta\sim N(0,\Delta+\Delta R_p\Delta).
\]
Its filtered field satisfies
\begin{equation}\label{cap:initial-frozen}
 R_f=(2-W+t+\alpha+(1-\beta)W)^{-1},\qquad a_f=R_fz,
 \qquad
 a_f\sim N(a_p,R_p-R_f)\quad\text{under }\mathsf G_h.
\end{equation}
In particular,
\[
 \mathbb E\|a_f\|^2=n-\operatorname{Tr}R_f
       =(1+o(1))n\sqrt{t+\alpha}.
\]
Under $\mathsf G_h$, the centered noise has covariance norm $O_{C_*}(t)$
and depends on the scalar $\alpha$ in a compact smooth family. Its shift
$h+\Delta a_p$ has normalized norm $O_{C_*}(t^{5/4})$. The scalar and
recipe nets in Proposition~\ref{cap:tap-expansion} cover this family and
all input directions. The exact mean above, the spectral trace bounds,
and a Gaussian quadratic deviation estimate give positive rates per $k_p$
for the fixed filtered-length tolerances.

For comparison, let $\mathsf T_h$ be the true terminal law conditional on
$\mathscr F_p^{\rm aux}$ and $h$:
\[
 X\sim\mu_{p,h},\qquad Z=h+\Delta X+\sqrt\Delta\,\xi,\qquad
 \xi\sim N(0,I)\ \text{independently of }X.
\]
The drift $h+\Delta X$ has normalized norm $O_{C_*}(t)$ for every spin.
The noise diagnostics hold on one event for all permitted shifts. The
low-band lower bound has a different quantifier: on a band with
$2-W\asymp t$ below $t$, condition first on $X$ and apply the Gaussian
small-ball bound to that fixed center, then average over $X$. This gives
failure $e^{-c nt^{3/2}}=e^{-cC_*^{3/2}k_p}$ with no union over centers.

The exact filtered Gaussian law in \eqref{cap:initial-frozen} is not the
law under $\mathsf T_h$. For those tests the spherical reference is the
canonical bridge conditioned on the prior spin norm. The tested norm-density
calculation gives a positive rate for the named relaxed quadratic tests.
The positive-moment transfer then attaches independent Gaussian noise to
each replica and uses the specified Gram-cell enlargement with the noise
marks fixed. Its field-uniform conclusion and parameter order are recorded
after the terminal regions are defined. The same interface is needed for
the actual finer frame at the starting cap.

The initial convolution uses one terminal potential for all input fields.
We construct it on three nested regions $D_0\subset D_1\subset D_2$,
determined by the parent data and the auxiliary frame. The region $D_0$
contains the predicted terminal fields, $D_1$ supplies the supporting
quadratics, and $D_2$ is the larger region on which the TAP Hessian is
controlled.

Fix $\eta_*>0$ at the choice of $C_*$, sufficiently small for the spectral
contraction and the Gaussian distance-tail comparison below. Choose a fixed
geometric tolerance $\theta$ sufficiently small compared with $\eta_*$.
The earlier finite choice of the annulus tolerance is made so that
$\sigma\ll\theta$. To define membership from a terminal field, let
$\Lambda=\{\mathrm{obs}\}\sqcup[-\sigma p^2,\sigma p^2]$ label the noise
covariances
\[
 \Sigma_{\mathrm{obs}}=\Delta,\qquad
 \Sigma_\alpha=\Delta+\Delta R_{p,\alpha}\Delta .
\]
Let $\mathcal G_W$ be the set of labelled noise vectors $(\lambda,g)$
passing the buffered diagnostics of Proposition~\ref{cap:tap-expansion}
for every allowed shift. Use closed relaxed versions of these tests and a
fixed Gaussian norm cutoff; their stricter versions retain the stated
probability bounds. The diagnostic families are continuous on their compact
coefficient sets, so $\mathcal G_W$ may be taken with a Borel graph and compact
sections. The low-band lower bound is a test of the resulting field, not a
condition imposed simultaneously on all shifts in $\mathcal G_W$.

Put $R_t=(2-W+t)^{-1}$ and $P_t^W=\mathbf1_{\{2-W\le t\}}$.
Choose shift radii $L_0^{\rm shift}<L_1^{\rm shift}<L_2^{\rm shift}$
inside the allowed range, with the first containing every drift
$h+\Delta X$ above and with fixed margins between the radii. Choose lower
thresholds $\ell_0>\ell_1>\ell_2>0$ below the small-ball threshold, and
$(\theta_0,\theta_1,\theta_2)=(\theta/4,\theta/2,\theta)$. Define
\begin{equation}\label{cap:terminal-regions}
\begin{split}
 D_j=\bigl\{z=g+v:\;&(\lambda,g)\in\mathcal G_W,\quad
 |v|_n\le L_j^{\rm shift}t,\\
 &|P_t^Wz|_n\ge\ell_jt^{5/4},\quad
 \bigl||R_tz|_n^2/\sqrt t-1\bigr|\le\theta_j\bigr\},
 \qquad j=0,1,2.
\end{split}
\end{equation}
Membership means that at least one labelled noise and shift pair satisfies
these conditions. All labels are included before $h$ is chosen. These are
nested parent-measurable closed sets: the defining graph has compact
sections, and its projection under $(g,v)\mapsto g+v$ is compact.
The geometric tolerances remain fixed throughout the argument.

The next lemma supplies the terminal tests at the fixed initialization
scale. Its filtered-length estimate has a rate available before the replica
order and upper scale are chosen. The noise diagnostics are then imposed at
that chosen scale; their rate need only be positive there.

\begin{lemma}[Terminal tests at the initialization scale]
\label{cap:marked-terminal-tests}
Fix the preceding geometric tolerances and a filtered-length tolerance
$\varepsilon_{\rm len}>0$. For a sufficiently small fixed $p=p_0$ and
then sufficiently small fixed $1-\beta>0$, with probability tending to one
in the capped data and auxiliary refinement, uniformly for
$h\in\mathcal A_p(\sigma)$,
\begin{equation}\label{cap:terminal-test-input}
 \mathsf T_h(D_0^c)\le e^{-a k_p},
\end{equation}
and
\begin{equation}\label{cap:terminal-filtered-test}
 \mathsf T_h\!\left\{
 \left|\|R_fZ\|^2-(n-\operatorname{Tr}R_f)\right|
       >\varepsilon_{\rm len}n\sqrt t\right\}\le e^{-a k_p}.
\end{equation}
Here $a>0$ may depend on $p_0$ and the fixed parameters. For either
exceptional event, the same conclusion holds, with constants $C_j,a_j>0$,
after inserting any fixed power of
$1+\|Z\|/\sqrt n+\|R_fZ\|/M_p$ in the expectation. A finite collection
of stricter filtered-length tolerances may be imposed simultaneously.
The same assertions hold for the actual refinement at this starting cap.
\end{lemma}

\begin{proof}
We first prove the filtered tests, keeping their rate uniform as the upper
scale decreases. Under the canonical joint law, sample
$X\sim N(a_p,R_p)$ and independent standard Gaussian $\xi$, and set
$Z=h+\Delta X+\sqrt\Delta\,\xi$. With $A=R_fZ$ and $T=R_p-R_f$,
Gaussian conditioning gives the joint representation
\begin{equation}\label{cap:terminal-joint-gaussian}
 A\sim N(a_p,T),\qquad X=A+\eta_f,\qquad
 \eta_f\sim N(0,R_f)\ \text{independently of }A.
\end{equation}
In particular $\mathbb E\|A\|^2=n-\operatorname{Tr}R_f$. The spectral
estimates, uniformly for $|\alpha|\le\sigma p^2$ and
$(1-\beta)/t$ sufficiently small, give
\[
 \|T\|\le Cp^{-2},\qquad
 \operatorname{Tr}T^2+a_p^{\mathsf T}Ta_p\le Cn/p.
\]
The Gaussian quadratic moment formula therefore implies, for each fixed
$\varepsilon>0$,
\begin{equation}\label{cap:canonical-terminal-length}
 \Pr\!\left\{\left|\|A\|^2-\mathbb E\|A\|^2\right|
             >\varepsilon n\sqrt t\right\}
 \le 2e^{-c_{\varepsilon,C_*}k_p}.
\end{equation}
Indeed its centered logarithmic moment is at most
$C\lambda^2 n/p$ for $|\lambda|\le cp^2$; choose a sufficiently small
fixed multiple of $p^2$ in each sign. These constants are independent of
sufficiently small $p$.

We must retain this rate after conditioning $\|X\|^2=n$. For a spectral
quadratic tilt $e^{\lambda A^{\mathsf T}HA+v^{\mathsf T}A}$, the normalized
tilted law of $X$ is Gaussian with covariance
\[
 R_f+T^{1/2}
   (I-2\lambda T^{1/2}HT^{1/2})^{-1}T^{1/2}.
\]
Choose the tilt with its bracket between fixed positive multiples of $I$.
This covariance is then comparable to $R_f+T=R_p$, regardless of the
tilted mean. A band of order $k_p$ eigenvalues remains comparable to
$p^{-2}$. The characteristic-function estimate
\eqref{cap:norm-characteristic} bounds its squared-norm density by
$Cp^2/\sqrt{k_p}=C\sqrt{p/n}$. The original saddle density is at least
$c\sqrt{p/n}$. Thus the density ratio in
\eqref{cap:conditioned-tail} is bounded by a constant, and
\eqref{cap:canonical-terminal-length} holds with a decreased positive
rate under the spherical terminal law. Independent marked replicas have
the corresponding product bound; equivalently one may use
\eqref{cap:joint-norm-density}.

This argument also covers the $R_tZ$ test defining $D_0$. Write
$R_tZ=(R_tR_f^{-1})A$. The spectral multiplier is bounded and differs
from $I$ by $O(|\alpha|/t+(1-\beta)/t)$, so its quadratic tilt has the same
precision margin. Its canonical mean square divided by $n\sqrt t$ is
$1+o(1)+O(\sigma/C_*)+O((1-\beta)/t)$. The trace estimate and the
resolvent identity give this comparison. Choose these errors below a
fixed fraction of $\theta_0$ and reserve fixed gaps between the tested
and relaxed thresholds. The spherical failure rate for both filtered
tests is consequently a fixed positive multiple of $k_p$ before $p_0$
is chosen.

We now transfer these concrete tests to cube posteriors by
\eqref{cap:tested-transfer}. Give every replica its own independent
standard Gaussian mark in spectral coordinates, and keep the marks fixed
in the Gram-cell comparison. A cell of the size used there changes a
spin representative by $O_l(n^{-1/2})$ in Euclidean norm. Hence it changes
$Z$ by $O_l(tn^{-1/2})$ and $R_fZ$ or $R_tZ$ by
$O_{l,C_*}(n^{-1/2})$. The corresponding changes in the normalized
filtered lengths tend to zero at the fixed scale. Norm inequalities in
both directions show that a failure with tolerance $\varepsilon$ implies
failure with tolerance $\varepsilon/2$ throughout the cell, even for unbounded
noise marks. This verifies the marked inclusion required after
\eqref{cap:tested-transfer}. Integrating the independent marks leaves
the same cube--sphere Gram factor. The small-overlap moment bound thus
has exponent $-lc k_p/2+\eta_l k_p+o_n(k_p)$; the discarded large-overlap
term is added after normalization as in that equation.

For uniformity in $h$, take a mesh of fixed resolution $\zeta$ in
$p^2M_p$ units on a slightly enlarged annulus. Its logarithmic cardinality
is at most $C_\zeta k_p$. Couple neighboring fields by using the same
spin and mark. The scaled saddle is Lipschitz, and the resolvent identity
gives
\[
 \|R_{f,h}Z_h-R_{f,h'}Z_{h'}\|
 \le C\zeta\bigl(M_p+\|R_{f,h'}Z_{h'}\|\bigr).
\]
Also $|\operatorname{Tr}R_{f,h}-\operatorname{Tr}R_{f,h'}|
\le C\zeta np$, so the centers of the squared-length tests respect the
same fixed gaps. The $R_t$ comparison is simpler, since its resolvent is fixed. Thus a
sufficiently fine fixed mesh respects the threshold gaps. On
$\|P_pX\|\le VM_p$, changing the spin posterior between these fields
costs at most $e^{C\zeta V k_p}$, including the partition normalizer;
use the static projected-mean bound on a prescribed ball containing the
annulus and the joining segment. The
complement has the uniform projected $V^6$ tail from
Proposition~\ref{cap:static-comparison}. Choose $V$ first to give the
required tail rate and then $\zeta$ to fit both the threshold and
likelihood budgets. In \eqref{cap:tested-transfer}, choose a posterior
exponent below $c/2$, the value error below the remaining margin, and
then a fixed replica order $l$ whose Haar exponent exceeds $C_\zeta$.
Only afterward decrease $p_0$ so that the Gram error $C_lnp^4$ and all
comparison errors fit this budget. The field net then proves both
filtered tests uniformly with a positive exponent per $k_p$.

It remains to impose the other conditions of $D_0$. The deterministic
shift $h+\Delta X$ is within the chosen shift radius for every spin.
Lemma~\ref{cap:boundary-noise} gives one conditional noise event for all
these shifts, with failure at most $e^{-c_tn}$ for typical $W$. For the
low-band lower bound, condition on each spin and use that lemma's
uniform-in-center small-ball estimate, then average. This bound is
uniform in the spin law and in $h$, without a field net or a union over
spins. At the chosen $p_0$, the minimum of these positive rates and the
filtered-test rate is a positive rate per $k_{p_0}$, proving
\eqref{cap:terminal-test-input}.

At this fixed scale, the conditional Gaussian observation and
$\|X\|=\sqrt n$ bound every fixed moment of
$1+\|Z\|/\sqrt n+\|R_fZ\|/M_p$, uniformly in $h$. Cauchy--Schwarz
with a doubled moment gives the weighted exceptional bounds with a
smaller exponent. The auxiliary refinement has marginal GOE law and is
independent of the true finer frame conditional on the capped data.
The event just proved is therefore measurable before that true frame is
revealed. Repeating the same argument gives the actual-refinement event;
it is imposed separately, without including it in
$\mathscr F_p^{\rm aux}$. A finite intersection proves the final assertion.
\end{proof}

The true terminal field now lies in $D_0$ with the weighted error bounds
needed for gradient transport. We next extend the TAP potential while
preserving its values and derivatives on a neighborhood of these fields.
A sufficiently small fixed $b_*>0$ gives a ball of radius
\begin{equation}\label{cap:agreement-buffer}
 b_*\sqrt n\,t^{5/4}
\end{equation}
inside $D_1$ around every point of $D_0$, and inside $D_2$ around every point
of $D_1$. Indeed, retain a witnessing $(\lambda,g)$ and assign the
displacement to $v$. A displacement of normalized length $b_*t^{5/4}$
changes the low-band norm by at most that amount and changes
$|R_tz|_n^2/\sqrt t$ by at most $Cb_*+Cb_*^2$. The fixed gaps between the
shift, lower-bound, and filtered-length thresholds give both buffers.

We also reserve an outer shift and low-band margin beyond $D_2$. Choose one
residual radius $\rho_{\mathrm{res}}$ inside these margins, with
$\rho_{\mathrm{res}}\ll t^{5/4}$. If
$e=F_{J_*,z}(y)$ and $|e|_n\le\rho_{\mathrm{res}}$, then
$F_{J_*,z+e}(y)=0$. The same noise witness remains admissible at $z+e$,
and the realized low-band bound is preserved by the triangle inequality.
The pre-stability part of Proposition~\ref{cap:tap-expansion} therefore
applies to every such approximate root. Its zero-field exclusion first
bounds $q$ above, the low-band estimate bounds it below by $c\sqrt t$,
and its scalar constraint, together with the filtered-length condition and
monotonicity, places it in the common band used below.
The derivative estimates then give the stability implication on this
whole residual region, with common parameters. Only at this point do we
apply the deterministic root lemma in Proposition~\ref{cap:classical-input}:
it gives a globally unique stable root for every $z\in D_2$.
The root is measurable in the data and field, for example as the limit of
measurably selected rational approximate roots; locally it is smooth by
the inverse function theorem.

The local curvature and the comparison of distant support points require
different arguments. To make the latter comparison without changing the
geometric regions, use the fixed interval
$I_t=[(1-2\theta)\sqrt t,(1+2\theta)\sqrt t]$ and define
\begin{equation}\label{cap:own-saddle-potential}
 R_q=(2+q^2-W)^{-1},\qquad
 F_q(z)=\tfrac12 z^{\mathsf T}R_qz+\frac n3q^3,
 \qquad
 \partial_qF_q(z)=q\bigl(nq-\|R_qz\|^2\bigr).
\end{equation}
The resolvents are positive on $I_t$ for the fixed boundary scale and large
$n$. The function $q-\|R_qz\|^2/n$ is strictly increasing. Its value at
$\sqrt t$ has absolute value at most $\theta\sqrt t$ on $D_2$, so it has a
unique zero $\widehat q(z)$ in the interior of $I_t$. Thus $\widehat q(z)$ is
the minimizer of $F_q(z)$ on this same interval for every $z\in D_2$.

The value calculation in Proposition~\ref{cap:tap-expansion} gives the exact
coefficient $\kappa_0=1/3$. Together with its mean expansion it implies,
uniformly on $D_2$, for a field-independent constant $c$,
\begin{equation}\label{cap:own-saddle-accuracy}
 \left|\Psi(z)-c-\min_{q\in I_t}F_q(z)\right|
       \le\epsilon nt^{3/2},\qquad
 \left\|\nabla\Psi(z)-R_{\widehat q(z)}z\right\|
       \le\epsilon\sqrt n\,t^{1/4}.
\end{equation}
Here $\epsilon$ is an approximation tolerance, separate from $\theta$ and
the two buffers. It can be decreased after those geometric quantities are
fixed. For clarity, the TAP resolvent initially has parameter
$r+r^{-1}=2+q^2/(1-q)$. Replacing it by $2+q^2$ changes its quadratic value
by $O(nq^4)$ and its mean by $O(\sqrt n\,q^{3/2})$, uniformly on the
coarse band $q\asymp\sqrt t$. The mean estimate and the definition of the
TAP overlap give $q-\|R_qz\|^2/n=o(q)$. Strict monotonicity then gives
$q-\widehat q(z)=o(q)$. On this fixed band the first two scalar derivatives
of $F_q/n$ have orders $q^2$ and $q$, so this replacement changes the value
by $o(nq^3)$ and the mean by $o(\sqrt n\,q^{1/2})$. These errors, and the
perturbations from the interaction coefficient and diagonal completion, give
\eqref{cap:own-saddle-accuracy}: choose the diagnostic accuracies, then the
boundary scale sufficiently small, then the coefficient perturbation
sufficiently small, and finally $n$ large. The noise diagnostic accuracy may improve, but assigning displacements to
the allowed shift still preserves those diagnostics. These choices do not
reduce the geometric tolerances or either fixed buffer.

By the resolvent identity and the fixed choice $\theta\ll\eta_*$,
$R_q\preceq R_t+\eta_*I/(4t)$ on $I_t$. The local TAP Hessian, whose leading
correction is negative semidefinite, therefore obeys the required upper
bound on $D_2$. Put
\begin{equation}\label{cap:initial-support-curvature}
 U=(2-W+t)^{-1}+\frac{\eta_*}{t}I,
 \qquad \|\Delta^{1/2}U\Delta^{1/2}\|\le1-b_{**}
\end{equation}
for a fixed $b_{**}>0$. The last inequality follows in the common spectral
basis: before the $\eta_*$ and $1-\beta$ perturbations,
$1-\Delta_i/(d_i+t)\ge c/C_*$ on both sides of the cap. The approximation
tolerance is also chosen small enough that $\nabla^2\Psi\preceq U$ on $D_2$.

For $z,z'\in D_1$ at distance less than
\eqref{cap:agreement-buffer}, their segment lies in $D_2$; integration of
this Hessian bound gives the upper-support inequality. For more distant
points put $d=z'-z$ and choose the minimizing parameter at $z$.
Since $\min_{q\in I_t}F_q(z')\le F_{\widehat q(z)}(z')$,
\eqref{cap:own-saddle-accuracy} gives
\[
 \Psi(z')-\Psi(z)-\nabla\Psi(z)^{\mathsf T}d
 \le\tfrac12d^{\mathsf T}R_{\widehat q(z)}d
       +2\epsilon nt^{3/2}
       +\epsilon\sqrt n\,t^{1/4}\|d\|.
\]
At these distances the two error terms are at most
$\epsilon(2b_*^{-2}+b_*^{-1})\|d\|^2/t$. Choose
$\epsilon\ll\eta_*\min(b_*^2,b_*)$ after the fixed buffers. The curvature
slack then bounds the right side by $d^{\mathsf T}Ud/2$. This is a
one-sided comparison using the saddle at the support point; it does not
require an arbitrarily accurate absolute comparison with a single frozen
quadratic on the whole geometric region.

Take the infimum of the upper supporting quadratics with curvature $U$
over $D_1$ and call the result $\overline\Psi$. The preceding inequalities
hold for every pair of support points. They show that the extension equals
$\Psi$ on $D_1$, hence on a full fixed neighborhood of $D_0$. The extension
has global upper Hessian $U$ in the distributional sense and at most
quadratic growth; its support gradients have at most linear growth. One may
mollify and pass to the limit at fixed $n$.

For later estimates we also need to transfer canonical tests to the heat
convolution of this extension. At each input field, restrict the supports
further to a stricter predicted subset $H_h$ of $D_0$ where $q/\sqrt{t+\alpha}$
and the filtered norm at $R_f$ have arbitrarily small fixed errors. This subset has canonical probability
tending uniformly to one. On it, $\Psi$ and its gradient agree with the frozen
canonical quadratic up to an arbitrarily small error per $k_p$ and per $M_p$,
respectively. The distance to this subset, in $\sqrt n\,t^{5/4}$ units, has
Gaussian upper tails with a positive rate per $nt^{3/2}$, uniformly as $p$
decreases. Indeed, Gaussian enlargement uses only
$\mathsf G_h(H_h)\ge1/2$ and the covariance bound $C_{C_*}t$;
the large-$n$ threshold for this mass bound may depend on the fixed $p$,
while the distance-tail rate does not. The support comparison bounds the additional logarithmic weight by
a small constant times $k_p$ plus
\[
 C(\eta_*+\sigma)nt^{3/2}
      \operatorname{dist}(z,\text{the subset})^2,
\]
where distance is in those scaled units. Choose $\eta_*$ below the
distance-tail rate. Integrating the distance tails, or applying H\"older's
inequality, transfers the required tests with a smaller positive exponent; the
same subset gives the normalization lower bound. For tests excluding $D_1$,
the tolerances and exponents can be fixed
independently of small $p$. Only then decrease $p$ and $1-\beta$ to control
the approximation errors.

For the lower Fisher estimate, localize before differentiating the extension.
Choose a cutoff equal to one on the stricter predicted subset in $D_0$ and
supported strictly inside $D_1$, where $\overline\Psi=\Psi$. This ensures
that the calculation uses only the original smooth potential, even though
$D_1$ need not be convex. The $D_0$-to-$D_1$ buffer
\eqref{cap:agreement-buffer} gives polynomial derivative bounds for a
distance cutoff. The larger domain $D_2$ is used for the local Hessian
segments in the support comparison. Integration by parts with the cutoff and the
score uses only the classical Hessian where the extension agrees locally.
Negative distributional Hessian mass elsewhere is never integrated. The
omitted terms have exponential probability and polynomial moments for the
fixed parameters.

\subsection{The initial radial deficit}

\begin{proposition}\label{cap:auxiliary-initialization}
For the preceding choices, with probability tending to one there is
an intrinsic auxiliary cap potential
\[
 \widetilde\phi_p(h)=
 \log\mathbb E_g\exp\overline\Psi(h+\sqrt\Delta\,g),
 \qquad h\in P_p,
\]
whose Hessian satisfies \eqref{cap:invariant} on the used annulus,
with slack.  Moreover, for a finite constant $C$ independent of $n$,
\begin{equation}\label{cap:initial-gradient-error}
 \sup_{h\in\mathcal A_p(\sigma)}
 \|\nabla\widetilde\phi_p(h)-P_p\nabla\phi_p(h)\|\le C.
\end{equation}
The constant may depend on all the fixed initialization parameters.
\end{proposition}

\begin{proof}
We first verify curvature. Write $\widetilde{\mathsf T}_h$ for the terminal
field law with density proportional to $e^{\overline\Psi(z)}$ relative to
$N(h,\Delta)$. All expectations in the curvature calculation below use this
law. An indicator $\mathbf1_{H_h}$ means that the integrand is integrated only
over $H_h$, including when it is defined only there. The canonical identities
in \eqref{cap:initial-frozen} enter through
the transferred tests, rather than as identities under
$\widetilde{\mathsf T}_h$.

On $H_h$, put $\nu=n/(4q)$; no value of $q$ or $\nu$ outside $H_h$ is used
in these retained integrals. The ideal form of the TAP response there is
\[
 K_f^0=R_f-\frac{R_fa_fa_f^{\mathsf T}R_f}
                    {\nu+a_f^{\mathsf T}R_fa_f}.
\]
Proposition~\ref{cap:tap-expansion} has the same two resolvent factors with
the root vector $y$. Replacing its resolvent by $R_f$ and
$y=r^{-1}a_f+o(M_p)$ by $a_f$ changes the response by an arbitrarily small
amount in $t^{-1}$ units; the scalar matching is detailed below. In the
ideal calculation, $R_f^{-1}=R_p^{-1}+\Delta$ and rank-one inversion give
the exact identity
\begin{equation}\label{cap:initial-rank-transport}
 \Delta^{-1}(\Delta^{-1}-K_f^0)^{-1}\Delta^{-1}-\Delta^{-1}
 =R_p-\frac{R_pa_fa_f^{\mathsf T}R_p}
              {\nu+a_f^{\mathsf T}R_pa_f}.
\end{equation}
The left side is the Brascamp--Lieb covariance bound after the
increment-covariance transformation. Thus the transported correction has
two $R_p$ factors, and its denominator contains $R_p$.

The transferred canonical tests give
\[
 \mathbb E[\mathbf1_{H_h}a_f]=a_p+o(M_p),\qquad
 \|a_p\|^2=(b+o(1))np,
\]
and
\[
 \begin{split}
 \mathbb E[\mathbf1_{H_h}a_f^{\mathsf T}R_pa_f]
 &\le a_p^{\mathsf T}R_pa_p+
       \operatorname{Tr}\bigl(R_p(R_p-R_f)\bigr)+o(n/p)\\
 &\le (3b+o(1))n/p.
 \end{split}
\]
Here $b=4/(3\pi)$, while
$\mathbb E[\mathbf1_{H_h}\nu]=O(n/(\sqrt{C_*}p))$.
On $H_p$, $R_p=(p^2+\alpha)^{-1}I$. For
$e_p=a_p/\|a_p\|$, Cauchy--Schwarz therefore bounds the retained expected
radial subtraction in $p^2$ Hessian units below by
\begin{equation}\label{cap:initial-radial-deficit}
 \frac{p^2}{(p^2+\alpha)^2}
 \frac{\bigl(\mathbb E[\mathbf1_{H_h}e_p^{\mathsf T}a_f]\bigr)^2}
      {\mathbb E[\mathbf1_{H_h}(\nu+a_f^{\mathsf T}R_pa_f)]}
 \ge \frac{1}{3+o(1)+O(C_*^{-1/2})}.
\end{equation}
The errors also tend to zero with $\sigma$. Taking $C_*$ large leaves a
strict margin above the radial constant $0.30$.

For each fixed unit transverse direction $v\perp a_p$, the retained second
moment $\mathbb E[\mathbf1_{H_h}(v^{\mathsf T}(a_f-a_p))^2]$ is $o(np)$,
with bounds uniform in that direction. The Fisher lower bound and the
Brascamp--Lieb upper bound therefore give arbitrarily small transverse error
and a bounded radial lower error. The contribution from $H_h$ in the localized
Fisher calculation is linear in the terminal correction and retains its denominator
$\nu+a_f^{\mathsf T}R_fa_f$; the upper bound uses the denominator in
\eqref{cap:initial-rank-transport}. The directional second-moment test
controls both on transverse directions.

The reference resolvent must be matched at the actual frozen parameter.
Put $s=t+\alpha$ and $q_0=\sqrt s$. On $H_h$, the
$R_f$ filtered-length test is within $\epsilon_{\rm len}q_0$ of its exact
mean $1-\operatorname{Tr}R_f/n$, and the latter is within
$\epsilon_{\rm tr}q_0$ of $q_0$. The TAP mean expansion gives
$|q-|R_qz|_n^2|\le\epsilon_{\rm TAP}q$. Since
\[
 \frac{d}{dv}\bigl(v-|R_vz|_n^2\bigr)
 =1+\frac{4v}{n}z^{\mathsf T}R_v^3z\ge1,
 \qquad
 R_f^{-1}=R_{q_0}^{-1}+(1-\beta)W,
\]
these tests imply
\[
 \frac{|q-q_0|}{q_0}
 \le C\left(\epsilon_{\rm len}+\epsilon_{\rm tr}
            +\epsilon_{\rm TAP}+\frac{1-\beta}{t}\right).
\]
Consequently the TAP parameter $q^2/(1-q)$ differs from $t+\alpha$ by
an arbitrarily small fraction of $p^2$ after the prescribed accuracy,
scale, and coefficient choices. This uses the exact $R_f$ test, so no fixed
$O(\sigma)$ error remains. The exceptional set is already covered by the
terminal transition bounds. The transverse error can therefore be made as
small as required at fixed $C_*$ and $\sigma$.
Positive-matrix Cauchy--Schwarz applied to the two matrix bounds gives the
mixed-block estimate in \eqref{cap:invariant}. The cutoffs described above
justify the lower bound, and \eqref{cap:initial-support-curvature} controls
the upper bound on the exceptional set.

The curvature invariant is now established for the auxiliary potential. We
next compare its gradient with the true cap mean in unnormalized Euclidean
norm. At the terminal field this is the uncentered consequence of the
companion's directional estimate. We transport that estimate backward through
the fixed precision interval. This gives a bounded gradient error; it does not
yet give a Hessian comparison.

At every terminal field in the agreement region the companion
fixed-$\beta<1$ posterior theorem gives
\begin{equation}\label{cap:accepted-mean-error}
 \|m_{\rm Gibbs}(z)-m(z)\|\le C.
\end{equation}
Precisely, use \cite[Section~7.2, proof of Proposition~7.2]{AcceptedGap}: its
estimate $|\mathbb E[G(X-t_*)^{\mathsf T}v]| \le C\|G\|_*$, before centering
$G$, with $G=1$ and $\|v\|\le1$ gives \eqref{cap:accepted-mean-error}. Its
assumptions are the fixed-length word diagnostics and the stable-field
implication through a fixed residual threshold. After fixing $\beta<1$,
restrict the permitted diagonal enlargement to a positive
$\eta<\min(\eta_{\rm TAP},\beta^{-1}-1)$, where $\eta_{\rm TAP}$ is the
enlargement supplied by the boundary test. The stability implication remains
valid on this smaller class, and $\beta^2(1+\eta)^2<1$ supplies the
companion's strict diagonal margin. Impose its fixed-length word diagnostics
with this choice. All constants are uniform in the terminal field and $n$
under these diagnostics. No assertion about the constant as $\beta\uparrow1$
is needed.

Transport \eqref{cap:accepted-mean-error} backward through the observation,
using Lemma~\ref{cap:error-transport} below. The strict upper curvature bound
\eqref{cap:initial-support-curvature} gives bounded amplification on this
fixed precision interval. Outside agreement, true observation probabilities
and their polynomially weighted moments are exponentially small by
Lemma~\ref{cap:marked-terminal-tests}. The extension gradients have at
most linear growth, so their contribution there tends to zero. This proves
\eqref{cap:initial-gradient-error}. Finally the conditional auxiliary
completion uses only the parent data and an independent Haar seed. Its
high-probability properties may be imposed before revealing the true finer
frame, which proves the asserted intrinsic measurability. \end{proof}

\subsection{Transporting and contracting gradient errors}

\begin{lemma}\label{cap:error-transport}\label{cap:gradient-transport}
Consider a deterministic increasing matrix precision $B_s$, with
derivative $G_s\succeq0$, and a terminal potential extended with a
strict semiconcavity margin relative to the total remaining
precision.  Let $f_s$ be its backward log heat convolution and
$K_s=\nabla^2 f_s$.  Along the true observation process, write
$m_s$ for the posterior mean and $e_s=m_s-\nabla f_s$.
For a starting direction $v$, let $v_s$ solve
\begin{equation}\label{cap:linear-flow}
 dv_s=dB_s\,K_s v_s=G_sK_sv_s\,ds.
\end{equation}
Then
\[
 e_{s_0}^{\mathsf T}v
      =\mathbb E[e_{s_1}^{\mathsf T}v_{s_1}\mid h_{s_0}],
\]
whenever the indicated polynomial moments are finite.  The same
identity holds by approximation for the extensions used here.
\end{lemma}

\begin{proof}
The true posterior mean has zero drift. Apply It\^o's formula to
the approximate gradient along the same observation, and use the backward
heat equation for $f_s$. Subtracting the two evolutions gives
$de_s=-K_sG_se_s\,ds+dM_s$.
The finite-variation equation \eqref{cap:linear-flow} cancels this
drift in $e_s^{\mathsf T}v_s$.  Localization and then the moment
bounds give the identity.  A global upper Hessian bound for the
terminal potential gives a global upper bound for the convolved
Hessian by Brascamp--Lieb.  For a constant precision rate $G$
this bounds amplification in the $G^{-1}$ metric.  On the cap
path the Euclidean estimate applies instead, because the
variation stays in the support of $P_u$, on which the rate is
scalar.  Mollification,
the same bound, and dominated convergence justify the limiting
extension.  In \eqref{cap:linear-flow}, $dB_s$ is a precision
increment, not a Brownian increment.
\end{proof}

\begin{proposition}\label{cap:gradient-contraction}
Suppose the parent auxiliary potential is intrinsic, satisfies the
curvature invariant with slack, and its true gradient discrepancy
is bounded by $D/p$ on its agreement annulus, where $D<\infty$
is fixed.  Put $c=\gamma p$.  Extend the parent potential as in
Proposition~\ref{cap:curvature-induction} and take its log heat
convolution to $c$.  For the parameter choices above, with
probability tending to one over the unused refinement,
\begin{equation}\label{cap:gradient-contraction-bound}
 \sup_{h\in\mathcal A_c(\sigma)}
 \|P_c\nabla\phi_c(h)-\nabla\widetilde\phi_c(h)\|
       \le\frac1c\bigl(\tfrac12D+o_n(1)\bigr).
\end{equation}
This conclusion is uniform in the child field.  For a fixed finite
number of comparison steps, the additive errors tend to zero
simultaneously.
\end{proposition}

\begin{proof}
Parameterize observation from $c$ to $p$ by $\log u$, so its rate
is $2u^2P_u$.  A variation initially in $P_c$ remains in $P_u$.
Freeze the starting parameter $\alpha$ and write
$a_u=R_uh_u$ along the canonical bridge.  These vectors form a
Gaussian martingale with independent increments and
\[
 \operatorname{Cov}(a_u-a_s)=R_s-R_u\qquad(c\le s\le u\le p).
\]
Norms and mixed products concentrate at their predictions, in particular
$a_s^{\mathsf T}a_u=\|a_s\|^2+o(M_p^2)$, with arbitrarily small fixed scaled
errors. The inferred parameter is within an arbitrarily small fraction of
$u^2$ of its frozen value. These assertions are uniform in $u$: use finite
time grids and Gaussian maximal bounds. In parent units the covariance of an
increment has operator norm at most $O_\gamma((\operatorname{Tr}P_p)^{-1})$
times its logarithmic time length and trace at most $O_\gamma(1)$ times that
length. Transfer of path tests to the true law uses its terminal density. The
exceptional-tail estimates are uniform conditionally on the starting field.
Outside the retained paths the global semiconcavity estimate still bounds flow
amplification.

If $h_u^{\rm lin}$ denotes the variation in \eqref{cap:linear-flow}, rescale
it to $v_u=(c^2+\alpha)h_u^{\rm lin}/(u^2+\alpha)$. The scalar reference
curvature cancels. By the intermediate estimates of
Proposition~\ref{cap:curvature-induction}, the remaining equation, up to a
small operator perturbation, is
\begin{equation}\label{cap:rank-flow}
 \frac{dv_u}{d\log u}
       =-2d_u\frac{a_ua_u^{\mathsf T}}{\|a_u\|^2}v_u,
       \qquad 0.26\le d_u\le C.
\end{equation}
The upper bound is absolute: use the variance-formula lower Hessian bound when
$u/p\le0.9$, and the Fisher bound closer to the parent. The perturbation is
bounded by $O_\gamma(\delta^{b'})$ for an absolute $b'>0$, plus errors chosen
arbitrarily small by the approximation tolerances. Variation of constants
absorbs it after these choices. The starting parameter tolerance contributes
no additional fixed scalar error, because the rescaling uses the actual frozen
parameter. All intermediate support conditions have slack; near $p$ use the
parent agreement domain itself.

Let $V_u:P_c\longrightarrow P_u$ be the pure flow
\eqref{cap:rank-flow}, with the identity inclusion at $c$.
It is a contraction.  Its terminal radial component satisfies
\begin{equation}\label{cap:radial-flow}
 \left\|\frac{a_p^{\mathsf T}V_p}{\|a_p\|}\right\|
       \le C\sqrt H\,\gamma^{1.02}+\tau,
       \qquad H=C_1\log(1/\gamma),
\end{equation}
where $\tau>0$ can be chosen arbitrarily small before $n$ is large. Removing
the scalar reference growth leaves a flow that contracts only in the current
radial direction. That direction changes along the bridge. To measure the
cumulative contraction, test the flow against the terminal radial vector and
use the bridge's mixed scalar products. The resulting scalar integral equation
has an integrating factor bounded by a fixed power of the scale ratio. We use
pathwise integration by parts, so no adaptedness claim is needed for that
factor.

To see the exponent, put $G_s=a_s^{\mathsf T}V_s$.
The mixed-product estimates in the integral equation give
\[
 G_s=a_s^{\mathsf T}V_c-
       \int_c^s2d_uG_u\,\frac{du}{u}
       +\text{an arbitrarily small scaled error}.
\]
Write $Q(p,u)=\exp(-\int_u^p2d_s\,ds/s)\le(u/p)^{0.52}$.
The integrating-factor formula is
\[
 G_p=Q(p,c)a_c^{\mathsf T}V_c+
       \int_c^pQ(p,u)\,d(a_u^{\mathsf T}V_c)
       +\text{an arbitrarily small scaled error}.
\]
The integral is interpreted by pathwise integration by parts; its integrand
can depend on future values of $d_s$, so an adapted stochastic-integral
interpretation is unnecessary. No differentiability of $d_s$ is needed. On a
logarithmic dyadic interval starting at $u$, increments and partial increments
of $P_ca_s$, divided by $\|a_p\|$, are at most $C\sqrt H\,\gamma^{3/2}p/u$
with exponential probability. This follows from the covariance formula and
$\operatorname{Tr}P_c\asymp nc^3$. Integration by parts against the
integrating factor gives a summable series with largest term $C\sqrt
H\,\gamma^{3/2}\gamma^{-0.48}$. The initial term is at most
$C\gamma^{1/2}\gamma^{0.52}$. This proves \eqref{cap:radial-flow}. Time and
mixed-product errors have at most polynomial amplification in $\gamma^{-1}$.

We next prove the transverse gain, including its conditional law. Let
$\mathcal F_p$ contain the capped matrix, the eigenvalues, and the entire
intrinsic auxiliary construction. Write $Q$ for the unused Haar map into
$P_p$. For a fixed child-coordinate field, sample the canonical path in
spectral coordinates independently of $Q$. Condition on these path coordinates
and on the terminal physical image $Q\xi_p=h_p$, where $\xi_p$ is the terminal
field in spectral coordinates. The terminal discrepancy, as a function of
$(\mathcal F_p,h_p)$, is now fixed. Its transverse part $e_p^\perp$ is
perpendicular to $h_p$, hence to $a_p$. The residual map from $\xi_p^\perp$ to
$h_p^\perp$ is exactly Haar. We do not condition on the intermediate physical
field images.

The elementary Haar projection bound, applied conditionally
as just specified, therefore gives
\[
 \|P_cQ^{\mathsf T}e_p^\perp\|
       \le C\sqrt H\,\gamma^{3/2}\|e_p^\perp\|,
 \qquad
 |(e_p^\perp)^{\mathsf T}a_u|/\|a_u\|
       \le C\sqrt H\,\gamma^{3/2}\|e_p^\perp\|
\]
on a fixed polynomially fine time grid, with failure at most $\exp(-A Hk_c)$
for a prescribed fixed $A$ after increasing the constant. Here $P_c$ in the
first expression denotes the child coordinate subspace before applying $Q$.
The bound uses $\operatorname{Tr}P_c/\operatorname{Tr}P_p \asymp\gamma^3$;
adding finitely many grid directions does not affect it once $n$ is large.
Gaussian modulus tests fill the gaps with any prescribed fixed error. Since
$\|V_u\|\le1$ and $d_u\le C$, the integral equation for $V_p$ yields
\begin{equation}\label{cap:transverse-flow}
 \|V_p^{\mathsf T}e_p^\perp\|
 \le\bigl[C\sqrt H\,\gamma^{3/2}
                (1+|\log\gamma|)+\tau\bigr]\|e_p^\perp\|.
\end{equation}
This step also permits the coefficients $d_u$ to depend on $Q$. The projection
event controls the vectors appearing in the integral equation, while their
coefficients use only the deterministic bounds already established. No
additional independence of those coefficients is required.

The residual Haar statement is used only under canonical
sampling.  To pass to the actual heat-tilted path, its
Radon--Nikodym derivative for each fixed $Q$ is the terminal
weight
\[
 \frac{\exp\{\phi_p(Q\xi_p)-q_p^{\rm fr}(\xi_p)\}}
 {\mathbb E_{\rm can}
    \exp\{\phi_p(Q\xi_p)-q_p^{\rm fr}(\xi_p)\}},
\]
where $q_p^{\rm fr}$ is the frozen quadratic. An additive calibration of the
numerator is immaterial. On an arbitrarily narrow fixed terminal shell, the
static comparison and the auxiliary gradient calibration bound this derivative
by $\exp(\eta k_p)$ for any prescribed $\eta>0$, and the same shell gives its
normalization lower bound. This normalization may depend on $Q$. We transfer
the canonical Haar failure estimate through the density ratio, rather than
asserting that the tilted law is Haar. If a finite grid or replica transfer
has $r$ terminal factors, choose $r\eta\ll H\gamma^3$. More explicitly, if
$\mathcal E_H$ is the high-rate Haar failure event on the retained terminal
shell, the canonical bound and the $r$ density factors give
\[
 \mathbb E_Q\,\mathsf T_{h,Q}^{(r)}(\mathcal E_H)
 \le \exp\{-[AH-r\eta\gamma^{-3}]k_c\}.
\]
Here $\mathsf T_{h,Q}^{(r)}$ is the joint actual path law for those
$r$ terminal factors at the fixed grid field.
Choose $b_{\rm path},b_Q>0$ with
$b_{\rm path}+b_Q<AH-r\eta\gamma^{-3}$ and
$b_Q>C_{\rm net}\log(1/\gamma)$, where the child-field net has logarithmic
size at most $C_{\rm net}\log(1/\gamma)k_c$. Markov's inequality then gives
\[
 \Pr_Q\!\left\{\mathsf T_{h,Q}^{(r)}(\mathcal E_H)
                 >e^{-b_{\rm path}k_c}\right\}
 \le e^{-b_Qk_c}.
\]
The second exponent pays for the field net. The first remains as a
conditional path bound and absorbs the polynomial discrepancy and flow
moments at the fixed comparison scales. Here is the likelihood comparison
needed to pass from the net to nearby fields. Let
$\|h-h'\|\le\rho c^2M_c$ and keep the later physical grid fields fixed.
Only the first Gaussian increment and the initial normalizer then change.
Its inverse on $P_c$ is bounded by a fixed multiple of $c^{-2}$ once the
first positive grid time is fixed. On bounded retained grid fields, the
Gaussian exponent comparison therefore gives
\[
 \left|\log\frac{d\mathsf T_h^{(r)}}{d\mathsf T_{h'}^{(r)}}\right|
       \le rC_\gamma\rho k_c.
\]
The normalizer obeys the same bound: differentiate the exact full-interval
convolution to obtain
$\nabla_{P_c}\phi_c=\Delta_{P_c}^{-1}(\mathbb E Y_{P_c}-h)$,
and use the uniform parent-field moment bound before integrating along
the field segment. The constant $C_\gamma$ may also depend on the fixed
first grid time and buffers.

The intrinsic terminal discrepancy is the same function of the fixed
terminal field in this comparison, so no derivative bound on it is
required. Discarding an arbitrarily short initial grid interval costs an
arbitrarily small operator error in the rank flow. The remaining
projection and mixed-product tests, including their normalizations by
nonzero annular lengths, have polynomial continuity in the fixed
parameters. Choose a net of inverse-polynomial resolution $\rho$ so
that these changes fit the reserved test slack and
$rC_\gamma\rho<b_{\rm path}/2$. A relaxed bad event at a neighboring
field is then bounded by its net-point probability times
$e^{rC_\gamma\rho k_c}$, retaining exponent at least
$b_{\rm path}/2$. This comparison is made before taking conditional
expectations. Its net still has logarithmic size
$C_{\rm net}\log(1/\gamma)k_c$, paid by the previously reserved
orientation exponent. Modulus tests,
mixed-product tests, and shell complements are controlled by the separate
uniform conditional tail estimates. Their failure exponents need not equal the
high Haar rate, and no union bound over the child net is applied to those
tails.

Apply Lemma~\ref{cap:error-transport}. The prefactor from undoing the
rescaling is $O(\gamma^{-2})$. Multiplying the resulting bound by $c$, the
radial and transverse terms in \eqref{cap:radial-flow} and
\eqref{cap:transverse-flow} are bounded by
\[
 C\sqrt H\left[\gamma^{0.02}
           +\gamma^{1/2}(1+|\log\gamma|)\right]D.
\]
Choose $\gamma$ so this is below $1/4$, and choose the operator and
fixed-tolerance errors so their contribution is below $D/4$. The exceptional
sets contribute $o_n(1)$ in these units for fixed $D$ and fixed scales. This
proves \eqref{cap:gradient-contraction-bound}. Agreement at each extension
boundary is enough: true paths leave it only with the stated negligible
probability. Integrating the bounded gradient discrepancy along annular paths
and recalibrating an additive constant gives an $o(k_c)$ value oscillation
error, since $c$ is fixed during this finite phase. The next curvature step
therefore retains all its macroscopic comparison hypotheses. \end{proof}

\subsection{Replacing the auxiliary potential and evaluating all scales}

After finitely many contractions the gradient discrepancy is small in the
scale units needed for the next convolution. This is the point at which a
Hessian comparison becomes available. The logarithmic density discrepancy is
Lipschitz on the parent shell; strong log-concavity turns that Lipschitz bound
into small thermal fluctuations. Changing density therefore perturbs the
directional first and second moments of the terminal field by a small amount.
The exact increment covariance formula then compares the two child Hessians.

\begin{lemma}\label{cap:remove-auxiliary}
Fix $\gamma$ and the curvature tolerances.  If the parent
gradient discrepancy is at most $D/p$ on the used annulus,
then one further convolution makes the true and auxiliary
output Hessians differ by $O_\gamma(D)+o_n(1)$ in child
$c^2$ units.  In particular, for $D$ sufficiently small,
the true output satisfies \eqref{cap:invariant} whenever
the auxiliary output satisfies it with the established
slack.
\end{lemma}

\begin{proof}
On the parent shell the difference of logarithmic
potentials has Lipschitz constant at most $CD/p$.
Annular paths have a bounded ratio of path length to
chord length, so this is also a Euclidean Lipschitz
bound up to an absolute factor.  Extend the difference
with the same bound, for example by the infimum formula
for Lipschitz extensions.  Under the strongly log-concave
extended auxiliary bridge, its centered fluctuations
are subgaussian with scale $C_\gamma D$.  Indeed the
bridge covariance scale is $O_\gamma(p^2)$ and its
strong convexity bound applies to every Lipschitz
function.  Linear functions of the terminal field in
thermal units have a bounded subgaussian scale as well.

If $F$ is the centered logarithmic discrepancy and
$L$ such a normalized linear function, Cauchy--Schwarz
and the subgaussian exponential moments give
\[
 \left|\frac{\mathbb E[e^F L^j]}{\mathbb E e^F}
                      -\mathbb E L^j\right|
       \le C_\gamma D,\qquad j=1,2,
\]
for $D$ below a fixed constant. Center the linear observable in each tested
direction before applying these inequalities. Their constants are uniform over
unit directions, so the resulting variance estimates control the full
covariance operator. The true bridge and the tilted extension agree
conditionally on the shell. The true bridge has the preceding complementary
moment bounds. For the discrepancy-tilted extension, H\"older's inequality
and the centered subgaussian bound on $F$ transfer the auxiliary complementary
moment bounds with a decreased positive exponent, including the required
second moments. The increment covariance
formula, with inverse increment $O_\gamma(p^{-2})$ on $P_c$, now proves the
asserted Hessian comparison. \end{proof}

Starting with Proposition~\ref{cap:auxiliary-initialization}, write its finite
error bound as $D_0/p_0$. Apply Proposition~\ref{cap:gradient-contraction} a
fixed number of times until $D$ is as small as
Lemma~\ref{cap:remove-auxiliary} requires, and perform the final replacement
step. The number of steps may be chosen after $p_0$, $\beta$, and $D_0$ are
fixed. There is no restriction on its size as a function of $p_0$: it is
finite and independent of $n$. The tolerances in the curvature induction and
in the flow estimates depend only on the fixed ratios and curvature
tolerances, not on this number of steps or on $D_0$. Raw bounded gradient
errors give vanishing macroscopic calibration errors at every one of these
fixed scales. Thereafter all cap potentials in the induction are the true
potentials. The resulting curvature conclusion depends only on the original
disorder: if it fails for a given $J$, the joint construction fails for every
auxiliary seed. Its disorder failure probability is therefore bounded by the
joint failure probability.

\begin{theorem}[Uniform curvature of the true cap potentials]
\label{cap:true-curvature}
Fix a finite field radius $L_0$. There is $C=C(L_0)<\infty$ such that
for every $\epsilon>0$ there are constants $p_{\rm top}>0$ and $s_0>0$
for which, with probability tending to one,
the following assertions hold simultaneously for every
\[
 p_{\min}(n)\le p\le p_{\rm top},\qquad
 p_{\min}(n)=\left(\frac{\sqrt{\log n}}{n}\right)^{1/3}.
\]
For all $h\in P_p$ with $|\alpha_p(h)|\le s_0p^2$,
\begin{equation}\label{cap:true-sharp-curvature}
 P_p\nabla^2\phi_p(h)P_p
       \preceq \frac{1+\epsilon}{p^2}P_p.
\end{equation}
For all $h\in P_p$ with $\|h\|\le L_0p^2M_p$,
\begin{equation}\label{cap:true-ball-curvature}
 P_p\nabla^2\phi_p(h)P_p\preceq\frac C{p^2}P_p.
\end{equation}
The constants are independent of $n,p,h$.  The field
radius and all coarse constants may be fixed before
the sharp curvature tolerances.  The lower endpoint
$p_{\min}(n)$ may be changed by a fixed factor.
\end{theorem}

\begin{proof}
Fix $L_0$ first and take the corresponding coarse constant from
Lemma~\ref{cap:frozen-evaluation-support}. Choose the construction
tolerances after this radius and the requested sharp error $\epsilon$;
the coarse constant is independent of $\epsilon$.
After the finite auxiliary phase, iterate
Proposition~\ref{cap:curvature-induction} on a geometric
list down to $np^3\asymp\sqrt{\log n}$.  Its
positive-rate failure bounds are summable on this list.
The static comparison and spectral estimates hold on
the same event.  This proves the invariant for all
retained true parent scales.

For \eqref{cap:true-sharp-curvature}, write the target scale as $p=xq$,
where $q$ is a true parent and $x\in[\gamma^2,\gamma]$. Choose $s_0$
small enough that the target frozen parameter lies strictly inside the
parent agreement range. Lemma~\ref{cap:deterministic-evaluation} applies
on this entire ratio interval. Its transverse and mixed errors can be
made arbitrarily small, while its radial correction is negative. Thus
$p^2(K_p-R_p^H)$ has arbitrarily small positive part. Choose the fixed
errors below the margin allowed by $\epsilon$, including the replacement
of $1/(p^2+\alpha)$ by $1/p^2$. This proves
\eqref{cap:true-sharp-curvature} without another random compression.

The sharp estimate concerns a thin parameter annulus. A prescribed bounded
ball may have a wider range of saddle parameters, so it requires a different
evaluation ratio. Choose the parent sufficiently coarser that this entire
range lies strictly inside the parent's agreement annulus. The child saddle
remains separated from its pole, and an upper support with smaller excess than
that separation gives the coarse ball bound. Each evaluation may use its own
support extension, since no induction is performed with its output.

For the ball bound again write $p=xq$, now with
$x\in[\gamma b_1,b_1]$, where
$b_1$ is supplied by Lemma~\ref{cap:frozen-evaluation-support} with
radius $L_0$. The lemma constructs supports around the frozen typical
length with excess below the child convolution margin, and proves the
complementary moment estimates for both the extended and true
transitions. Equation~\eqref{cap:ball-evaluation-bound} therefore gives
\eqref{cap:true-ball-curvature}, with a coarse constant depending only on
$L_0$. The construction is uniform in the indicated fixed ratio interval;
it uses no radial gain and need not agree with the extensions used for
other targets.

Every target scale below a sufficiently small fixed $p_{\rm top}$ has a parent
in each indicated ratio interval. The two evaluation lemmas are deterministic
and uniform under the parent hypotheses. Thus the assertions include all
intermediate scales, not just the geometric list. Fixed-factor enlargement of
the lowest list interval covers the stated bottom endpoint. \end{proof}
